% ==========================================
% SECCIÓN 2: ANÁLISIS DEL PROBLEMA
% ==========================================

\section{ANÁLISIS DEL PROBLEMA}

\subsection{Antecedentes Institucionales y Contexto Operativo}

La Unidad de Gestión de Doctorados de la Universidad Central del Ecuador (UCE) constituye la dependencia técnica de la administración central facultada para orientar, registrar, supervisar y auditar el trayecto académico-administrativo de los docentes titulares pertenecientes a las diversas facultades que se encuentran cursando programas de posgrado doctoral (Ph.D.) en instituciones de educación superior nacionales e internacionales, bajo esquemas de becas institucionales o licencias con remuneración. El seguimiento riguroso de este universo docente abarca el control continuo de múltiples variables contractuales, entre las que destacan la filiación del docente, la universidad de destino, los plazos de vigencia de las becas, la emisión y aprobación de adendas o prórrogas autorizadas, la verificación de los hitos de graduación y el cómputo estricto del periodo de devengamiento exigido legalmente tras la consecución del título doctoral.

Históricamente, ante la ausencia de un ecosistema informático diseñado de manera específica para la complejidad de estos procesos, la Unidad ha recurrido al uso de archivos de hojas de cálculo de Microsoft Excel almacenados de forma local en los equipos de trabajo del personal. Cada modificación en el estatus contractual de los docentes becarios, así como el recálculo de los plazos de devengamiento, debe ser registrado manualmente por el personal administrativo miembro de la Unidad, integrado por el Director de la Unidad, la Analista de Procesos Universitarios y el Analista Técnico. Si bien esta modalidad ofimática sirvió como un paliativo operacional en sus etapas iniciales, el sostenido crecimiento en la plantilla de docentes en formación doctoral ha sobrepasado los límites de eficiencia y confiabilidad de los mecanismos manuales.

\subsection{Planteamiento y Diagnóstico del Problema}

La gestión de expedientes doctorales a través de hojas de cálculo aisladas ha generado cuellos de botella operativos que comprometen la gobernanza de la información y la precisión administrativa institucional. En primer lugar, se evidencia un fenómeno agudo de fragmentación de datos y duplicidad de versiones, al no contar con un repositorio centralizado de datos que actúe como fuente única de verdad en tiempo real. La actualización independiente de archivos locales por parte de los analistas de procesos y del director conduce a discrepancias en las fechas registradas de inicio y término de becas, generando inconsistencias entre las distintas facultades y la administración central.

En segundo lugar, el cómputo de los plazos temporales —que incluye la duración de la beca, la aplicación de adendas reglamentarias y el cálculo del tiempo obligatorio de devengamiento— descansa en fórmulas manuales propensas al error humano. La falta de validaciones de integridad referencial y de un sistema automatizado de trazabilidad o auditoría expone a la institución a vicios de extemporaneidad en la notificación de vencimientos de contratos. Esta vulnerabilidad se incrementa al no contar con respaldos automatizados en la nube o en la infraestructura del servidor central, lo que representa un riesgo inminente de pérdida irreversible de datos ante eventuales fallas mecánicas de hardware o ciberataques.

En tercer lugar, el modelo operacional vigente restringe severamente la autonomía y comodidad de los docentes becarios al privarlos de un portal web de consulta integral y autogestión. Para conocer el estado actualizado de su expediente, verificar qué documentos tienen debidamente registrados y qué requisitos o recaudos les falta entregar, los docentes se ven obligados a realizar visitas presenciales recurrentes a la oficina de la Unidad o enviar insistentes correos electrónicos. Asimismo, ante la carencia de un sistema de reservas y agendamiento de citas en línea en tiempo real, aquellos becarios que requieren orientación personalizada con el Director o los Analistas deben acudir de manera espontánea, generando aglomeraciones, interrupciones laborales continuas y demoras en la atención. Esta dinámica, sumada a la ausencia de alertas automáticas y de un generador digital de certificados de estado de doctorado, satura al personal administrativo con consultas repetitivas que se solucionarían eficientemente mediante un portal de autoservicio estudiantil dotado de agendamiento de citas en tiempo real y módulo de solicitudes.

\subsection{Análisis Crítico y Prognosis}

El análisis crítico del problema demuestra que la causa estructural radica en la falta de un ecosistema informático desacoplado compuesto por un panel administrativo que reemplace las hojas de cálculo para la gestión de doctorados y estudiantes, un portal estudiantil interactivo con visualización clara de documentación (entregada vs. faltante) y reservas de citas en tiempo real, y una API REST centralizada que articule la comunicación segura entre ambas capas. El uso continuado de herramientas no especializadas degrada la calidad del servicio administrativo, sobrecarga al personal de la Unidad y genera malestar en los docentes becarios por la falta de canales modernos y transparentes de información.

En cuanto a la prognosis, de no implementarse una plataforma web modular automatizada, la sobrecarga del personal administrativo derivará en un colapso operacional inevitable. Esto se traducirá en el incumplimiento de los términos legales establecidos en la Ley Orgánica de Educación Superior (LOES) y el Reglamento de Régimen Académico UCE \cite{loes}, \cite{reglamento_uce}, la generación de contingencias legales o financieras por errores en el control de expedientes y plazos, y un incremento en la insatisfacción de los docentes becarios ante la opacidad, la necesidad de traslados físicos innecesarios y la demora en la atención de sus trámites.

\subsection{Sistematización del Problema (Preguntas de Investigación)}

Con el propósito de abordar sistemáticamente la problemática planteada, se estructuran las siguientes preguntas de investigación:

\begin{description}[font=\bfseries\color{UceBlue}]
    \item[Pregunta Principal:] ¿De qué manera el diseño e implementación de la plataforma web modular Sistema de Gestión de Doctorados (SGD), compuesta por un portal estudiantil con visualización de estado de expedientes y reservas de citas en tiempo real, un panel administrativo centralizado y una API REST desacoplada, optimizará la gestión de los programas de doctorado y la atención a los docentes becarios de la Universidad Central del Ecuador?
    \item[Subpregunta 1:] ¿Cuáles son las entidades, relaciones y reglas de integridad requeridas para diseñar una base de datos relacional PostgreSQL con Prisma ORM que consolide y normalice los expedientes doctorales, requisitos documentales, solicitudes y el calendario de citas, eliminando el uso de hojas de cálculo de Excel?
    \item[Subpregunta 2:] ¿Cómo estructurar un panel administrativo web en Next.js 16 con React 19 y control RBAC que permita al personal de la Unidad gestionar integralmente los programas de doctorado, monitorear el avance de los estudiantes, administrar la agenda de citas y procesar trámites de forma ágil?
    \item[Subpregunta 3:] ¿Cómo desarrollar un portal web estudiantil en Astro y React que proporcione a los docentes becarios visualización en tiempo real de su expediente (documentos entregados y faltantes), un sistema de reservas de citas en línea con el personal de la Unidad, emisión de certificados y alertas preventivas?
    \item[Subpregunta 4:] ¿De qué forma se deben construir los servicios API REST centralizados en Express.js bajo contratos desacoplados y compilar manuales interactivos en Fumadocs para garantizar una comunicación fluida entre interfaces y orientar una futura migración institucional?
\end{description}

\subsection{Idea a Defender (Hipótesis de Trabajo)}

El diseño e implementación de una demostración tecnológica del Sistema de Gestión de Doctorados (SGD) basada en una arquitectura web modular desacoplada —conformada por un panel administrativo en Next.js 16 para la gestión integral de programas doctorales sin hojas de cálculo, un portal estudiantil en Astro con visualización del estado del expediente (documentos entregados y pendientes) y sistema de reservas de citas en tiempo real, interconectados mediante una API REST en Express.js con PostgreSQL y Prisma ORM sobre un banco sintético ficticio basado en formatos institucionales— optimizará la eficiencia administrativa de la Unidad al centralizar el control de expedientes y citas, y proporcionará autonomía informativa a los docentes becarios reduciendo significativamente las visitas presenciales no planificadas.
